\chapter{Os números reais}
Números reais são objetos matemáticos que representam, intuitivamente, medidas.
Uma forma usual de visualizar o conjunto $\mathbb R$ dos números reais é pensar neste como o conjunto dos pontos de uma reta.
\subimport{reais}{figura}

A Teoria dos Números reais e do Cálculo Diferencial e Integral moderna faz uso da linguagem de conjuntos.
A linguagem da Teoria dos Conjuntos pode ser formalizada a partir da Teoria dos Conjuntos de Zermelo-Fraenkel com o Axioma da Escolha.
Tanto a Lógica quanto a Teoria dos Conjuntos são disciplinas profundas com uma rica vida própria.
Neste texto, não daremos uma introdução formal a tais disciplinas, nos restringindo a apresentar a linguagem básica necessária para tornar este texto autocontido.

\subimport{}{reaisConjuntos}
\subimport{}{reaisDefinicao}
\subimport{}{reaisOrdem}
